\chapter{Родитель curvature moment map
\texorpdfstring{$\Delta$}{Delta}-стабилизатора}
\label{ch:v10-h15-r2-delta-moment-map-parent}

Нейтральный rank-one $\Delta$-фон задаёт две traceless moment maps
\begin{equation}
 \mu_R=\operatorname{diag}\left(\frac12,-\frac12\right),
 \qquad
 \mu_4=\operatorname{diag}\left(-\frac14,-\frac14,-\frac14,\frac34\right).
\end{equation}
Их trace-нормы равны $1/2$ и $3/4$. Канонически нормированная комбинация
на $\Sigma$ есть
\begin{equation}
 \boxed{Q=6Y=6\mu_R-4\operatorname{ad}(\mu_4)},
\end{equation}
и имеет прежний спектр $(3,-3,7,1,-1,-7,3,-3)$. Тем самым отношение
Cartan-каналов и коэффициент $2$ при смешанном $TB$ действительно
фиксируются одним moment map, а не двумя portal-весами.

Однако положительная curvature-норма сама по себе даёт Hessian
\begin{equation}
 Q^2=\operatorname{diag}(9,9,49,1,1,49,9,9).
\end{equation}
Она имеет ранг $8$ и делает $R_2\oplus\overline R_2$ самым тяжёлым, а не
лёгким сектором. Поэтому нужный знак нельзя получить объявлением
$-\|Q\Sigma\|^2$ фундаментальной нормой.

Mapping-cone подсказывает другой механизм. Введём общий auxiliary-канал с
полным квадратичным Hessian
\begin{equation}
 \mathcal H_{\rm aux}=
 \begin{pmatrix}
 49I&Q\\ Q&I
 \end{pmatrix}.
\end{equation}
Он конгруэнтен положительной форме
\begin{equation}
 \operatorname{diag}(49I-Q^2,I)
\end{equation}
через обратимую triangular matrix. Поэтому полный parent
положительно полуопределён, имеет ранг $14$ и нуллитет $2$. После
устранения auxiliary-поля его Schur complement точно равен
\begin{equation}
 \boxed{49I-Q^2=G_Y}.
\end{equation}
Отрицательный $Q^2$ возникает как стандартный subtractive Schur term, а не
как отрицательная норма полного действия.

Это закрывает условную архитектуру, но не происхождение. В унаследованной
fixed-point auxiliary algebra cross-блок между $\Sigma$ и auxiliary
сектором равен нулю. Без него устранение даёт только universal $49I$, а не
gap. Нужно предъявить восьмимерный типизированный образ $Q\Sigma$ внутри
36-мерного self-adjoint fixed-point пространства mapping cone.

\begin{theorem}[Условный Schur-parent гиперзарядового gap]
Нормированный $\Delta$ moment map единственно задаёт оператор $Q=6Y$.
Положительный полный Hessian $\bigl(\begin{smallmatrix}49I&Q\\Q&I\end{smallmatrix}\bigr)$
имеет Schur complement $G_Y=49I-Q^2$, поэтому условно порождает точный
$R_2$-gap без отрицательной фундаментальной нормы.
\end{theorem}

\begin{nogocriterion}[Формальный auxiliary не является унаследованным]
Нельзя считать Schur-parent физическим, пока карта
$\Sigma\mapsto Q\Sigma$ не вложена в конкретный fixed-point auxiliary
модуль старого mapping cone с правильными reality, grading и trace metric.
Нулевая текущая cross-карта не может быть заменена абстрактным $I_8$.
\end{nogocriterion}

\begin{protocol}\begin{enumerate}
\item moment-map нормы $\mu_R/\mu_4$: \textbf{$1/2,3/4$};
\item нормированная комбинация: \textbf{$6\mu_R-4\operatorname{ad}\mu_4$};
\item positive moment-map Hessian rank: \textbf{$8$};
\item positive norm ordering: \textbf{неверный знак};
\item full Schur-parent rank/nullity: \textbf{$14/2$};
\item full inertia: \textbf{$(0_-,2_0,14_+)$};
\item Schur complement: \textbf{$G_Y$};
\item inherited auxiliary coupling rank: \textbf{$0$};
\item conditional/physical origin: \textbf{$3/4$ и $2/4$};
\item следующий гейт: \textbf{typed embedding в fixed-point auxiliary channel}.
\end{enumerate}\end{protocol}

Машинный сертификат сохранён в
\path{s2t_v10_particle_wrinkle_dislocation_callias_equal_charge_twist_m4_h15_r2_delta_stabilizer_moment_map_curvature_parent_origin_gate_results.json}.